\section{Evaluation}
\label{sec:eval}

We evaluate \sys by asking:
\begin{enumerate}[\bfseries {Q}1]
\item Does \sys record exactly the over-budget frames? (\autoref{sec:eval-accuracy})
\item Does the partition attribute an injected stall to the right cause? (\autoref{sec:eval-accuracy})
\item Does chain resolution separate kernel stalls from waiting on work? (\autoref{sec:eval-chain})
\item Does frame-scoped accounting capture hitches made of many small stalls? (\autoref{sec:eval-cuts})
\item Does the cause gate reject hitches that are not kernel-observable? (\autoref{sec:eval-gate})
\item Do the records explain the mechanisms behind reported stutters, and
  what can they not say? (\autoref{sec:eval-cases})
\item When hitches are rare, which tools record them, attribute them, and
  name the culprit? (\autoref{sec:eval-catch})
\item What does always-on tracing cost, and how does output scale? (\autoref{sec:eval-overhead})
\item How common are kernel-observable hitches in Linux games? (\autoref{sec:eval-prevalence})
\end{enumerate}

\paragraph{Setup} The results in this draft come from a development virtual
machine: Hyper-V on an Intel Core i7-12700H host, 10 vCPUs, 7\,GB of memory,
Ubuntu 24.04.5 with Linux 7.0 (\texttt{PREEMPT\_LAZY}; Ubuntu's 7.0.0-31 to
7.0.0-42 kernels over the course of this work, and 7.0.0-42 for the \bench
run below), and
Mesa 25.2.8 with lavapipe~\cite{lavapipe}, the only Vulkan device the VM
exposes. The VM has no PMU. Its timing is not representative of a gaming
machine, so we use it only for questions whose answer does not depend on
absolute timing (Q1--Q7), and report Q8 from it as a pilot. Bare-metal
overhead, real titles and Q9 require a GPU and are planned on
\tbd{[evaluation machine: AMD Ryzen 7 PRO 8700GE with Radeon 780M,
Deck-class limits imposed with cgroups]}.

\subsection{Frame Capture and Attribution Accuracy}
\label{sec:eval-accuracy}

\paragraph{Benchmark} \bench is a synthetic game loop. Its root thread burns a
calibrated slice of CPU each frame, hands a job to a worker and waits for it
on a futex, then presents: the frame marker is called at present entry, the
thread sleeps to the frame deadline, and the present-return marker is called
after. The loop paces to 14.6\,ms so that healthy frames stay under a
16.7\,ms budget. Every twentieth frame it runs one injector that stalls the
frame through a single kernel mechanism, and it logs every frame's id,
duration, injection and expected bucket as JSON. We run each of
\hbnClasses\ injector classes and one uninjected baseline for eight seconds
with \sys attached for five, and join records to ground truth on frame id.
Frames that straddle the start or end of tracing are excluded.

\begin{table*}[t]
\centering
\footnotesize
\resizebox{\textwidth}{!}{\setlength{\tabcolsep}{2.6pt}
\begin{tabular}{@{}llccccccccc@{}}
\toprule
 & & \textbf{Expected} & \textbf{Recorded /} & \textbf{Largest bucket} & \textbf{Top group} & \textbf{Longest stall /} & \multicolumn{2}{c}{\textbf{Share gate}} & \multicolumn{2}{c}{\textbf{Resolved gate}} \\
\cmidrule(lr){8-9}\cmidrule(l){10-11}
\textbf{Class} & \textbf{Injection} & \textbf{bucket} & \textbf{injected} & \textbf{= expected} & \textbf{= expected} & \textbf{excess} & \textbf{share} & \textbf{pass $\geq$0.5} & \textbf{share} & \textbf{pass $\geq$0.5} \\
\midrule
\hbTableRows
\bottomrule
\end{tabular}
}
\caption{Attribution on \bench (medians over injected frames). \emph{Top
group}: the record's first-ranked stall group has the expected cause; on-CPU
classes have no group (--). \emph{Longest stall / excess}: the root's single
longest stall interval as a fraction of the frame's overrun, the most a
per-stall threshold could see. \emph{Share}: the root's kernel-observable
time (all buckets but \bucket{oncpu} and \bucket{block\_present}) over the
overrun, as the prototype's in-kernel gate computes it. \emph{Resolved gate}: the same
with time spent waiting on a running task subtracted, computed offline from
the records. Generated by \texttt{plot/src/hitchbench\_results.py}.}
\label{tab:bench}
\end{table*}

\paragraph{Q1: exactly the over-budget frames} Across the \hbnRuns\ runs, \bench
logged \hbtotFrames\ fully traced frames, of which \hbtotOver\ exceeded the
budget. \sys emitted exactly one record for each of them and none for any
frame within budget (\hbtotMiss\ missed, \hbtotExtra\ extra). Its frame
times agree with \bench's own to within \hbmaxDtUs\,$\mu$s. In the
uninjected baseline, \hbbaselineOver\ of \hbbaselineFrames\ frames
exceeded the budget, and \sys emitted \hbbaselineRecords\ records.

\paragraph{Q2: the right cause} \autoref{tab:bench} and
\autoref{fig:breakdown} show the result per class. Every injected frame
was recorded, and the largest bucket of the root's partition named the
injected mechanism in \hbtotMatch\ of \hbtotInj. The two classes that fall
short are informative rather than wrong. In \texttt{fault}, the injector
touches a burst of fresh pages; \sys carves a median \hbfaultoncpufault\,ms
out of on-CPU time as \bucket{oncpu\_fault}, but the loop that touches the
pages runs \hbfaultoncpu\,ms of user code of its own, so plain
\bucket{oncpu} stays larger. In \texttt{io}, the cold read spends a median
\hbioblockio\,ms in I/O wait but also \hbiooncpu\,ms on-CPU copying from the
page cache, and in \hbioToppresent\ of \hbioCaught\ frames the present wait
that follows exceeds the I/O on this VM's fast virtual disk. In both cases
the mechanism is present in the record at its real size; ``largest bucket''
is too blunt a criterion. \todo{Adopt the per-level
diagnosis criterion (injected bucket $\geq X$\% of excess above the
normal-frame baseline) and use \texttt{O\_DIRECT} for the I/O injector.}

\paragraph{Which stall the record leads with} A record ranks the root's stall
groups (\autoref{sec:groups}) by total time and explains each through its
chain. Over the \hbtotRank\ injected frames of the classes whose cause is an
off-CPU wait (the on-CPU classes \texttt{fault} and \texttt{cpu\_spike} have
no group), the first group had the injected cause in \hbtotTop. The misses
are the \texttt{io} frames above, where the present wait's group outweighs
the I/O's. Replaying on the same records the rule \sys used before stall
groups, which explained the root's longest single blocked interval, gives
\hbtotOld: it fails every \texttt{preempt} frame, because a preemption is not
a blocked interval, and every \texttt{thousand\_cuts} frame, because no
single sleep is long.

\begin{figure}[t]
\centering
\pgfplotstableread{plot/data/hitchbench_breakdown.dat}\breakdown
\begin{tikzpicture}
\begin{axis}[
  xbar stacked,
  width=0.88\columnwidth, height=6.9cm,
  bar width=6.5pt,
  y dir=reverse,
  ytick=data,
  yticklabels from table={\breakdown}{class},
  yticklabel style={font=\ttfamily\scriptsize},
  xticklabel style={font=\scriptsize},
  xlabel={Mean time per injected frame (ms)},
  xlabel style={font=\scriptsize},
  xmin=0, xmax=45,
  enlarge y limits=0.06,
  set layers,
  legend cell align=left,
  legend style={font=\scriptsize, at={(0.5,1.02)}, anchor=south,
    legend columns=5, draw=none, /tikz/every even column/.append style={column sep=4pt}},
  legend image code/.code={\draw[#1, draw=none] (0cm,-0.08cm) rectangle (0.25cm,0.1cm);},
  axis x line*=bottom, axis y line*=left,
  xmajorgrids, grid style={black!12},
]
\addplot[fill=cOncpu, draw=none] table[x=oncpu, y expr=\coordindex]{\breakdown};
\addplot[fill=cOncpuK, draw=none] table[x=oncpukern, y expr=\coordindex]{\breakdown};
\addplot[fill=cSched, draw=none] table[x=sched, y expr=\coordindex]{\breakdown};
\addplot[fill=cFutex, draw=none] table[x=futex, y expr=\coordindex]{\breakdown};
\addplot[fill=cPoll, draw=none] table[x=poll, y expr=\coordindex]{\breakdown};
\addplot[fill=cIO, draw=none] table[x=io, y expr=\coordindex]{\breakdown};
\addplot[fill=cTimer, draw=none] table[x=timer, y expr=\coordindex]{\breakdown};
\addplot[fill=cPresent, draw=black!30] table[x=present, y expr=\coordindex]{\breakdown};
\addplot[fill=cOther, draw=none] table[x=other, y expr=\coordindex]{\breakdown};
\legend{oncpu, oncpu kernel, runnable+rq, futex, poll, io, timer, present, other}
\draw[on layer=axis foreground, black, dashed, thick] (axis cs:16.667,-0.6) -- (axis cs:16.667,11.6);
\end{axis}
\end{tikzpicture}

\caption{Composition of the root's frame for each injector class (mean over
injected frames). Each bar sums to the frame time; the dashed line is the
16.7\,ms budget.}
\label{fig:breakdown}
\end{figure}

\subsection{Chain Resolution}
\label{sec:eval-chain}

\paragraph{Q3} Two classes produce the same symptom on the root, a futex wait of
about 20\,ms on the worker, for opposite reasons. In
\texttt{worker\_block} the worker is itself blocked, reading a pipe that a
feeder thread writes after a 20\,ms sleep. In \texttt{worker\_cpu} the worker
simply computes for 20\,ms. By bucket the two are indistinguishable
(\autoref{tab:bench}). By resolution they are opposites: for
\texttt{worker\_block}, a median \hbworkerblockInherited\,ms of the wait is
inherited and \hbworkerblockWaitOncpu\,ms is waiting on a running task; for
\texttt{worker\_cpu}, \hbworkercpuInherited\,ms and \hbworkercpuWaitOncpu\,ms.
The chain for \texttt{worker\_block} is
\hbworkerblockChain\ in \hbworkerblockChainN\ frames: the worker, the feeder
the worker was waiting on, and the timer interrupt that ended the feeder's
sleep. For \texttt{poll}, where the root waits in \texttt{epoll\_wait} for
the feeder directly, \hbpollInherited\,ms is inherited through
\hbpollChain.

This is the distinction a per-cause total cannot make and the reason \sys
resolves waits online: a futex wait on a stalled thread is a kernel stall
one hop away, while a futex wait on a busy thread is the game's own CPU
work, out of scope, and should not count as a kernel stall.

\texttt{server\_queued} has the same symptom again, for a third reason:
the thread the root waits on loses its CPU without being preempted. The root
waits on a condition variable for a reply from a server thread that runs at
the lowest priority and is asleep when three CPU hogs take its CPU; woken,
it waits about 20\,ms in the run queue behind them. It is never preempted, so there is no preemptor to name. Its
hop records that it was kept runnable, with \hbserverqueuedInherited\,ms of
the root's wait inherited, and names the task that held its CPU when it
finally ran, flagged \emph{queued}: a hog in \hbserverqueuedHolder\ frames.
The one miss is the first injected frame after \sys attached. The server had
gone to sleep before the attach, so \sys holds no off-CPU interval for it,
and the root's wait resolves as waiting on a running task. Only the last
holder is named; another task that held the CPU earlier in the same wait is
not.

\subsection{Hitches Made of Many Small Stalls}
\label{sec:eval-cuts}

\paragraph{Q4} The \texttt{thousand\_cuts} injector makes the root sleep
about forty times for 0.5\,ms each. Every injected frame was recorded with
the timer bucket largest, a median of \hbthousandcutsNstalls\ stalls in the
frame, and a longest single stall covering only \hbthousandcutsWorstFrac\
of the overrun, against \hbsingleFracLo--\hbsingleFracHi\ for the classes
that inject one wait. The sleeps form one stall group, a median of
\hbthousandcutsTopGroupCount\ waits, ranked first in
\hbthousandcutsTop\ frames. A per-event
filter at any threshold above 0.5\,ms emits nothing for these frames by
construction, and one below it emits a comparable stream for every normal
frame. \todo{Run \texttt{perf record --off-cpu --off-cpu-thresh} and the
ftrace per-stall synthetic event on the same schedule and report capture
rate vs.\ threshold (M2), and sweep stall size and count.}

\subsection{Cause Gating}
\label{sec:eval-gate}

\paragraph{Q5} Two classes are negative controls, over budget for reasons that
are not kernel stalls: \texttt{cpu\_spike} burns 25\,ms of CPU on the root,
and \texttt{present\_long} lengthens the present wait. The prototype's gate
at $\sigma=0.5$ rejects both in every frame (median share
\hbcpuspikeShare\ and \hbpresentlongShare) and passes every positive class.
It also passes \texttt{worker\_cpu} (share \hbworkercpuShare), because the
gate counts the root's futex wait whether or not the thread it waited for
was stalled. Subtracting the resolved \textit{wait\_oncpu} time, which \sys
already computes, drops \texttt{worker\_cpu} to \hbworkercpuRShare\ and
rejects it in every frame (\hbworkercpuRGated\ pass), while every positive
class still passes in every frame but one: the first \texttt{server\_queued}
frame after attach, whose wait resolves as waiting on a running task
(\autoref{sec:eval-chain}; \autoref{tab:bench}, last column).
\todo{Move the resolved gate into the BPF program and re-run.}

\paragraph{vkcube} We also ran \sys on \texttt{vkcube-wayland} under lavapipe,
which rasterizes inside the present call. With the present bracket and no
cause gate, 98 of 210 frames (46.7\%) exceeded 16.7\,ms; in their records
96.8\% of the time was \bucket{block\_present}, with \texttt{gnome-shell}
named as the waker in 97 of 98, and a median of only 12\% of the overrun was
kernel-observable. With $\sigma=0.5$, 25 of 479 frames (5.2\%) were emitted.
Without the bracket, the same waits appear as \bucket{block\_poll} (86\% of
recorded time) in every one of 223 records: the bracket is what keeps
display pacing from being reported as a stall. The frames that survive the
gate have a median overrun of only 147\,$\mu$s, sub-millisecond jitter that
passes a ratio test easily because the denominator is tiny.
\todo{Add a minimum-overrun floor to the gate.}

\subsection{Case Studies}
\label{sec:eval-cases}

\paragraph{Q6} \bench injects each cause into its own frame loop. To test the
records on causes that come from outside the game, we reproduce the
mechanisms behind documented Linux gaming stutters around one unmodified,
game-shaped Vulkan program. It has a main thread (frame limiter, simulation,
job submission, upload, present) and three workers that draw the frame in
software, runs at a 60\,fps cap on four CPUs, and has no injectors; a hitch
is a frame over 20.8\,ms. Something else on the machine supplies each cause:
a background program that wakes periodically and burns CPU, another program
reading from a bandwidth-limited disk, a cgroup memory limit, a shim that
raises the frame limiter's timer slack, a server process the game calls 30
times a frame, a tool that confines the game's threads to one core, or a
program that leaves free memory in pieces. Ground truth comes from those
programs' own logs and from the game's per-frame profiler, never from \sys.
Each run lasts 45 to 60\,s, and every figure is computed one record at a
time against that ground truth. These runs reproduce mechanisms, not the
reported incidents: they show whether a record is right when the cause is
known.

\begin{table*}[t]
\centering
\footnotesize
\setlength{\tabcolsep}{3.5pt}
\begin{tabular}{@{}l>{\raggedright\arraybackslash}p{3.7cm}cccc>{\raggedright\arraybackslash}p{3.4cm}@{}}
\toprule
 & & & & \multicolumn{2}{c}{\textbf{Ranked first}} & \\
\cmidrule(lr){5-6}
\textbf{Case} & \textbf{Cause} & \textbf{Explained} & \textbf{Amount} & \textbf{by total} & \textbf{by stall time} & \textbf{What the record cannot say} \\
\midrule
1 \texttt{bgcpu} & a background program takes the game's cores & 99\% & 100\% & \multicolumn{2}{c}{first (one long wait)} & which preemptor took longest \\
2 \texttt{iomain}, \texttt{iothread} & an asset read from storage & 100\% & 91--94\% & 66--71\% & 81--88\% & why the disk was slow \\
2b \texttt{iodisk} & another program reads from the same throttled disk & 100\% & 99\% & 94\% & 100\% & who else was using the disk \\
3 \texttt{mem}, \texttt{memthp} & walking memory that was swapped out & 99--100\% & 95--96\% & 33--60\% & 99--100\% & what caused the memory pressure \\
4 \texttt{slack} & the frame limiter's sleep ends late & -- & 99\% & 0\% & 0\% & that the sleep was late \\
5 \texttt{ipc} & a server the game calls cannot get its CPU & 100\% & 97--105\% & 8--83\% & 98--100\% & who held the CPU before the last holder \\
6 \texttt{cpuspike} & CPU work in the game (negative control) & ``running'' 100\% & 99\% & \multicolumn{2}{c}{no stall blamed in 91\%} & what the CPU was running \\
7 \texttt{squeeze} & the game's threads confined to one core & 100\% & 96\% & \multicolumn{2}{c}{100\%} & which CPU; that the others were idle \\
8 \texttt{alloc} & new memory while memory is full & 97\% & 94\% & never & never & what was pushed out, and why \\
9 \texttt{compact} & huge pages while free memory is in pieces & 100\% & 78\% & never & never & that the kernel fell back to small pages \\
\bottomrule
\end{tabular}

\caption{Case studies (VM). \emph{Explained}: over-budget frames affected by
the cause whose record names the mechanism. \emph{Amount}: the record's time
for the mechanism over the game's or the cause's own measure of it, median.
\emph{Ranked first}: the mechanism's stall group is the record's first by
total, or by \emph{stall time}, a ranking replayed offline that leaves out
waits on threads that were running and timer sleeps; on-CPU kernel time is
never a group. Ranges cover a case's variants and, for case 5, a run before
and after the queued holder. Case 4's sleep is measured, not explained.}
\label{tab:cases}
\end{table*}

\paragraph{Records are right when the cause is known} In every case but one,
the record of nearly every affected frame names the mechanism, mostly at 91
to 105\% of the size measured by the game or by the cause's own log
(\autoref{tab:cases}). No noise frame's record names a background program,
another disk user or a swap-in. Where a mechanism is always present at a
smaller size, noise frames show it smaller: every server call costs a
wake-up (case 5), the game's threads always contend a little (case 7, where
how many cores' worth ran in the frame, 1.0 against 1.9, tells the two
apart), and an allocation just under the ground truth's threshold reclaims
too (cases 8 and 9). The negative control is answered as it should be: in
every spike frame the main thread was running about 12\,ms more than usual,
and in 91\% no stall was worth a quarter of the overrun.

\paragraph{The cases found what \bench did not} Three parts of the design exist
because a case failed. In case 1 the chain inherited only the waker's last
off-CPU interval and credited a median 40\% of a preempted worker's delay;
with the off-CPU history (\autoref{sec:chain}) it credits 100\% and names the
background program as the preemptor in 90\% of burst frames and in none of
the noise frames. In case 8, reclaim inside a page fault was counted twice,
as reclaim and as fault time. In case 5 the server was never preempted, so
no record could say who held its CPU; with the queued holder, records name
the background program in 163 of 167 burst frames, and in the other four a
kernel worker that ran just before the server. Stall groups matter most in
case 3: a walk faults 64 pages back from swap, none of them waiting longer
than about 0.7\,ms. Summed, the swap-in is the largest stall in every frame,
but the longest single wait was a swap-in in only 8 to 15\% of them. The
group's count equals the kernel's major-fault count for the thread, frame by
frame, in 99 to 100\% of frames.

\paragraph{Ranking by total is not enough} A frame's normal waits are large: the
main thread waits about 10\,ms for its workers and sleeps about 6\,ms in its
limiter every frame. A short load's I/O group therefore ranks first by
total in only 66 to 71\% of load frames, and a swap-in's in 33 to 60\%. The
offline ranking by stall time puts the right group first in 81 to 100\%. We
have not made it the record's: it would rank a late limiter sleep last, and
case 4 shows that a late sleep is sometimes the cause. \sys measures that
sleep exactly, 9.5\,ms against a usual 6, but cannot tell a sleep that
overshot from one that was asked to be long.

\paragraph{What a record cannot say} The last column of \autoref{tab:cases} is
a list of what \sys would have to record next, and \autoref{sec:limitations}
collects it. Faults, reclaim and compaction (cases 8 and 9) are measured but
never ranked, because on-CPU kernel time forms no group.

\subsection{Catching Rare Hitches}
\label{sec:eval-catch}

\paragraph{Q7} When a game hitches rarely, which tools' output shows each hitch,
and which shows its cause? We wrote down the protocol and its scoring rules
before the first run; three changes after short check runs, and one after
the measured run, are recorded with their reasons in the repository. Five of
the case studies are made rare, one injection every 10\,s for 120\,s:
\texttt{bgcpu}, \texttt{iothread} (the main thread waits for a streaming
thread reading an 8\,MB asset), \texttt{ipc}, \texttt{mem}, and
\texttt{cpuspike} as the negative control. Each tool runs alone, once per
case, in shuffled order and as its documentation says, with its own
processes on CPUs the game does not use. An \emph{injected hitch} is a frame
that the ground truth marks and the game's log shows over budget; a
\emph{noise hitch} is any other over-budget frame.

The tools are \sys; per-frame delay accounting~\cite{delayacct}, read for
every game thread at each frame mark; an ftrace flight
recorder~\cite{ftracehist} that keeps \texttt{sched\_switch} and
\texttt{sched\_waking} in an overwriting ring and snapshots it on a slow
frame; \texttt{perf sched record}; bcc's
\texttt{offwaketime}~\cite{gregg2019bpf}; and \texttt{chaingraph}, our
earlier wake-chain profiler. After the measured run we added two more, with
their rules written down before their own runs: \texttt{perf record
--off-cpu --off-cpu-thresh 1}~\cite{perf-offcpu-thresh}, built with the BPF
skeletons it needs, which emits a sample for each blocked interval of at
least 1\,ms; and Perfetto~\cite{perfetto-triggers} as a flight recorder, its
default 32\,MB ring of \texttt{sched\_switch}, \texttt{sched\_waking} and
the frame mark copied to a file by a \texttt{CLONE\_SNAPSHOT} trigger.
Perfetto cannot see frames, so a watcher of the layer's frame marks fires
the trigger when the time from one mark to the next is over budget, as a
game or a system daemon would. Outputs are scored mechanically, per injected
hitch, at three levels: \emph{recorded} (the output covers the frame),
\emph{attributed} (it puts the frame's extra time on the injected cause
without further analysis), and \emph{culprit} (it names the background
program, or for the read the device interrupt). ftrace, Perfetto and
\texttt{perf sched} keep raw events, which hold the cause whenever they hold the frame but
say nothing until someone analyzes them, so they are scored for recording
and for its cost in output. \texttt{offwaketime} and \texttt{chaingraph}
print one summary per run, so they are scored on the rank of the heaviest
summary line that carries the cause; \emph{visible} means among the ten
heaviest.

\begin{table*}[t]
\centering
\footnotesize
\setlength{\tabcolsep}{3.5pt}
\begin{tabular}{@{}lccccccc@{}}
\toprule
 & \textbf{\sys} & \textbf{delay acct.} & \textbf{perf off-cpu} & \textbf{ftrace} & \textbf{Perfetto} & \textbf{perf sched} & \textbf{offwaketime} \\
\textbf{Case} & \multicolumn{3}{c}{\emph{recorded / attributed / culprit}} & \multicolumn{3}{c}{\emph{recorded}} & \emph{rank of the cause} \\
\midrule
\texttt{bgcpu} & 8/8/8 of 8 & 11/11/-- of 11 & 11/0/0 of 11 & 8 of 8 & 11 of 11 & 10 of 10 & 39th \\
\texttt{iothread} & 12/12/12 of 12 & 12/12/-- of 12 & 12/10/0 of 12 & 11 of 12 & 12 of 12 & 12 of 12 & 22nd \\
\texttt{ipc} & 33/33/\textbf{0} of 33 & 28/0/-- of 30 & 33/0/0 of 33 & 29 of 33 & 33 of 33 & 32 of 32 & absent \\
\quad with queued holder & 29/29/29 of 29 & & & & & & \\
\texttt{mem} & 12/11 of 12 & 9/0 of 9 & 12/0 of 12 & 12 of 12 & 12 of 12 & 12 of 12 & 20th \\
\texttt{cpuspike} & 12/12 of 12 & 12/12 of 12 & 12/0 of 12 & 12 of 12 & 12 of 12 & 12 of 12 & -- \\
\midrule
MB per injected hitch & 0.08--0.10 & $<$0.01 & 1.3--3.7 & 15--42 & 45--89 & 25--33 & 0.01--0.02$^\ast$ \\
\bottomrule
\end{tabular}

\caption{Rare hitches (VM; one injection every 10\,s, one 120\,s run per
case and tool). Counts differ between tools because a burst makes a hitch
only when it lands in a frame without slack: same mechanism, not the same
frames. \texttt{mem} and \texttt{cpuspike} have no culprit. The \texttt{ipc}
row ``with queued holder'' reruns \sys alone after the change of
\autoref{sec:chain}; perf \texttt{--off-cpu} and Perfetto ran later, each in
a run of its own. $^\ast$The whole summary, divided by the injected
hitches.}
\label{tab:catch}
\end{table*}

\paragraph{Only \sys attributes and names} \autoref{tab:catch} shows the
result. \sys attributed 76 of 77 injected hitches; delay accounting
attributed 35 of 74. It has no category for a wait on another process
(\texttt{ipc}), and it books a swap-in wait as block I/O on the faulting
thread, 8.7 to 31\,ms a frame against about 0.15\,ms of swap-in delay, which
does not tell a swap-in from a file read (\texttt{mem}). perf
\texttt{--off-cpu} attributed 10 of 80. Its samples carry the blocked
thread's user stack and the interval's length, with no kernel stack and no
waker: they show the read in \texttt{iothread} (in the two misses the read
blocked in waits under 1\,ms, which perf only sums at the end of the run),
but preemption is never sampled, a wait on the server or on a swapped-out
page says nothing of why it was long, and the main thread's usual wait for
its workers looks like any stall (\texttt{cpuspike}). Only \sys names a
culprit: the background program in 8 of 8 \texttt{bgcpu} hitches and the
disk's interrupt in 12 of 12 \texttt{iothread} hitches. In \texttt{ipc} the
measured run named none, because the server was queued rather than
preempted; this result led to the queued holder, and a rerun of \sys alone
names the background program in 29 of 29 hitches and in none of 88 noise
hitches. Across the measured run's 255 noise hitches in the three cases
with a culprit, \sys never named one. Its one miss is in \texttt{mem}: 42\,ms
of the frame were the game writing pages to swap inside reclaim, inside a
swap-in, and the kernel stack was cut at 32 frames before the page fault,
so the rule's ``the stack is a page fault'' did not match.

\paragraph{Summaries lose rare hitches} In \texttt{offwaketime} and
\texttt{chaingraph} the cause was never among the ten heaviest lines, and
its share of the summary rounded to 0.0\% in every case. In \texttt{bgcpu}
and \texttt{ipc}, where the background program preempts rather than wakes,
it does not appear in either summary at all.

\paragraph{Raw recorders hold the frames, at a price} \texttt{perf sched}
and Perfetto recorded every injected hitch and the ftrace flight recorder
all but 5 of 77, at 155 to 900 times \sys's output per hitch (ftrace
155--454, \texttt{perf sched} 293--420, Perfetto 502--903, case by case).
Four of the ftrace recorder's misses were slow frames 22 to 57\,ms after
another slow frame. A snapshot swaps the ring with a spare buffer rather
than copying it, and the reader must save and clear the spare before the
next swap; a second slow frame before that leaves a gap. Perfetto's
\texttt{CLONE\_SNAPSHOT} copies the ring, and every one of its 479 triggers
left a snapshot. But each snapshot is the whole ring: at about 150\,KB/s of
scheduler events the 32\,MB ring does not fill in a 120\,s run, so each one
holds everything since tracing began, 45 to 89\,MB per injected hitch, and
in a longer session each would be the full 32\,MB. None of these tools says
anything about the cause until someone analyzes its output.

\paragraph{Limits} One rarity, one run per case and tool, the VM's fast disk and
no GPU stalls, and runs too short to fill Perfetto's ring. A sweep of
rarities would show how quickly the summary tools' ranks fall; we have not
run it.

\subsection{Overhead and Output Volume}
\label{sec:eval-overhead}

\paragraph{Q8} We measured the cost of always-on tracing as a pilot in the VM.
The tools are those of \autoref{sec:eval-catch} except delay accounting
(below), and \sys also without its interval log. perf \texttt{--off-cpu} and
Perfetto were measured later, each in a pilot of its own with \sys
alongside; the VM's speed drifts between pilots, so each is compared with
\sys in its own run. The workloads are the
case-study game at its 60\,fps cap and uncapped, the same game calling a
server 30 times a frame, and \texttt{perf bench sched pipe}, run inside the traced process
to measure the cost per context switch to the game and in another process
to measure what tracing the game costs everything else. Tools run in
shuffled order, with a run without any tool before every second tool run,
in three rounds of 15\,s. Each tool run is compared with its neighbouring
runs without a tool (median ratio, bootstrap 95\% interval), and slowdown
is taken from the workload's own log, never from the tool being measured.
In-kernel BPF time comes from separate runs with
\texttt{kernel.bpf\_stats\_enabled}.

\begin{table*}[t]
\centering
\footnotesize
\setlength{\tabcolsep}{5pt}
\begin{tabular}{@{}lcccccc@{}}
\toprule
 & \multicolumn{2}{c}{\textbf{Added ns per switch}} & \multicolumn{2}{c}{\textbf{BPF time, \% of a core}} & \textbf{Slow frames caught} & \textbf{Output} \\
\cmidrule(lr){2-3}\cmidrule(lr){4-5}
\textbf{Tool} & \textbf{traced process} & \textbf{other process} & \textbf{game} & \textbf{pipe} & \textbf{game / ipc} & \textbf{MB per hour} \\
\midrule
\sys & 479 & 409 & 2.10 & 7.50 & 27/27, 49/49 & 28--75 \\
\sys, no interval log & 569 & 413 & 2.14 & 7.40 & 34/34, 56/56 & 20--39 \\
\texttt{offwaketime} & 881 & 322 & 1.13 & 11.81 & no per-frame output & 0.08--0.11 MB in all \\
\texttt{chaingraph} & 430 & 18 & 2.17 & 13.87 & no per-frame output & 0.11--0.36 MB in all \\
\texttt{perf sched} & 325 & 73 & -- & -- & all & 9,646--25,215 \\
ftrace flight recorder & 134 & 168 & -- & -- & 10/21, 44/60 & 1,766--5,154 \\
\midrule
\multicolumn{7}{@{}l}{\emph{Separate runs, each with \sys alongside:}} \\
\sys & 800 & 780 & 2.41 & 8.08 & 29/28, 40/39 & 31--57 \\
perf \texttt{--off-cpu} & \textbf{2,463} & 834 & n/m & n/m & no per-frame output & 1,375--2,068 \\
\sys & 419 & 505 & 1.84 & 7.43 & 23/23, 39/39 & 22--61 \\
Perfetto flight recorder & $\approx$0 & 195 & -- & -- & 23/22, 65/65 & 2,557--15,766$^\dagger$ \\
\bottomrule
\end{tabular}

\caption{Overhead pilot (VM, three rounds of 15\,s per tool). Added time per
context switch is half the added time per \texttt{perf bench sched pipe}
round trip, against a baseline of about 4.2\,$\mu$s per switch; differences
of a few hundred nanoseconds are within this VM's noise. Slow frames caught
are over the game's own count inside the measured window. The ftrace and
Perfetto recorders and \texttt{perf sched} use no BPF; n/m: not measured.
$^\dagger$In 15\,s runs, where Perfetto's ring is far from full; a long
session would write more.}
\label{tab:overhead}
\end{table*}

\paragraph{Frame time} On the game workloads, the median frame time and the
frame rate stayed within $\pm$1\% of the runs without a tool for every tool
but one. At game rates, about 10,000 context switches a second, no other
tool's cost is visible above this VM's noise. perf \texttt{--off-cpu}, which
takes the user stack at every switch-out and keeps perf's usual on-CPU
sampling, cost the uncapped game 5.5\% of its frame rate (95\% interval
1.5 to 9.3\%).

\paragraph{\sys is not cheaper} Per context switch it is among the more
expensive tools (\autoref{tab:overhead}): about 480\,ns on the traced
process, where only \texttt{offwaketime} costs more, and perf
\texttt{--off-cpu}, about three times \sys in its own run; and among the
most on other processes, about 410\,ns against \texttt{chaingraph}'s 18 and
\texttt{perf sched}'s 73, level with perf \texttt{--off-cpu} in its run.
Perfetto costs almost nothing per switch, like the ftrace recorder, since
the kernel only writes each event into a ring; its work is done in user
space, where reading and encoding the events took 1.0 to 2.3\% of a core,
about what \sys's BPF programs take in the kernel (1.8\% in the same run). The reason is deliberate: it keeps a history for
every task on the machine, because a chain can run through any process
(\autoref{sec:limitations}). Its in-kernel time on the game is level with
\texttt{chaingraph}'s and about twice \texttt{offwaketime}'s; its switch
hook takes about 600\,ns per call there and its wakeup hook about 700\,ns.
The queued holder adds a few comparisons per switch and writes only after a
run-queue wait of at least 100\,$\mu$s; rerunning the pipe workloads after it
(three runs) gave 70 to 880\,ns per switch on another process across three
variants that run the same hook, so no change could be resolved.

\paragraph{What the cost buys} \sys's output is a few hundred times smaller than
\texttt{perf sched}'s, the only other tool here that keeps per-frame detail
of everything, and it recorded every slow frame, where the ftrace flight
recorder caught 10 of 21 and 44 of 60. That recorder misses frames whose
extra time was inside present: its frame runs from the return of one present
to the next frame mark, because its histogram trigger pairs two different
events and cannot pair present entry with itself. Perfetto, whose watcher
sees the whole frame, caught every slow frame too, at 2.6 to 15.8\,GB per
hour in these short runs. A record is 10,344 bytes
plus 72 per logged interval, so at one hitch per second \sys writes about
40\,MB per hour, and nothing for frames within budget. The summary tools
cost about the same as \sys but give one summary per run.

\paragraph{What we claim} \sys's overhead is in the range of existing scheduler
tracers, with no measurable effect on frame time in this VM. For that cost it
produces a record of causes for every slow frame at a few hundred times less
output than full tracing and misses none, so it can be left on for hours of
play. It does not have lower CPU overhead.

\paragraph{Not yet measured} Bare metal, with ten rounds of 30\,s; a cheaper path
for tasks outside the game, tracked fully only once they touch its wake
chain; \texttt{perf sched}'s writes interfering with a game that streams from
the same disk; and delay accounting, whose pilot numbers we discard: the
kernel accounts a thread only if accounting was on when the thread was
created, and the harness turned it on after the game had started.
\todo{Bare-metal run with the full preset; controlled hitch rates of 0.01,
0.1 and 1 per second; snapshot drain time and dropped snapshots.}

\subsection{Prevalence in Linux Games}
\label{sec:eval-prevalence}

\paragraph{Q9} \todo{Not yet measured. Plan (M0): \tbd{5} titles (Proton and
native) $\times$ 3 conditions (unconstrained; Deck-class cgroup limits;
Deck-class plus a background download and a cold shader cache), $\geq$45
minutes each, $\sigma=0$ so every over-budget frame is recorded. Report the
distribution of kernel-observable share of the overrun, resolved as in
\autoref{sec:eval-gate}; the breakdown by bucket; and the split between
game-initiated waits and external causes. This number has not been
published for Linux games and is the paper's motivating measurement either
way it comes out.}
